\section{The exam and how to use this guide}

\begin{itemize}
  \item \textbf{Format:} 3 hours, \textbf{40 multiple-choice questions}, 1 point each. A wrong answer and a blank answer both give 0, so \textbf{always answer every question}.
  \item \textbf{No aids allowed} (25H, 23H and 22H resit all say ``Legal aids: None'').
  \item \textbf{Weight:} the exam is 40\% of the final grade (the group project is 60\%).
  \item \textbf{When several options seem right}, the exam rule is: \emph{pick the one that fits best}. The questions often say ``What will be the best option that covers all possibilities''.
\end{itemize}

\subsection*{What gets asked (approximate question counts)}
\begin{center}
\begin{tabular}{lccc}
\toprule
Topic & 25H & 23H & 22H resit \\ \midrule
Python basics (types, lists, dicts, scope, try, files) & 10 & 15 & 12 \\
NumPy (indexing, axis, shape, broadcasting, sort) & 7 & 6 & 9 \\
Pandas (Series/DataFrame, merge, groupby, cleaning, I/O) & 12 & 6 & 11 \\
Time series (date\_range, shift, rolling, resample, tz) & 3 & 2 & 4 \\
Plotting (Matplotlib/Seaborn) & 2 & 7 & 1 \\
Machine learning theory & 6 & 4 & 3 \\ \bottomrule
\end{tabular}
\end{center}

\subsection*{How each section works}
Each section explains the concept from zero and then gives:
\begin{itemize}
  \item a \colorbox{ideabg}{\textbf{Key idea}} box: the one thing to understand;
  \item a {\color{exambar}\textbf{Asked in past exams}} box: real past questions, each with its answer;
  \item a \colorbox{trapbg}{\textbf{Traps}} box: the wrong options the exam uses to trick you.
\end{itemize}
Section~\ref{sec:tricks} at the end has the answer-key patterns and a last-minute checklist.
